
\documentclass{article}
\usepackage{colortbl}
\usepackage{makecell}
\usepackage{multirow}
\usepackage{supertabular}

\begin{document}

\newcounter{utterance}

\centering \large Interaction Transcript for game `codenames', experiment `risk\_high', episode 9 with reflexion{-}Qwen3.5{-}9B.
\vspace{24pt}

{ \footnotesize  \setcounter{utterance}{1}
\setlength{\tabcolsep}{0pt}
\begin{supertabular}{c@{$\;$}|p{.15\linewidth}@{}p{.15\linewidth}p{.15\linewidth}p{.15\linewidth}p{.15\linewidth}p{.15\linewidth}}
    \# & $\;$A & \multicolumn{4}{c}{Game Master} & $\;\:$B\\
    \hline

    \theutterance \stepcounter{utterance}  
    & & & \multicolumn{2}{p{0.3\linewidth}}{
        \cellcolor[rgb]{0.9,0.9,0.9}{
            \makecell[{{p{\linewidth}}}]{
                \texttt{\tiny{[GM$|$GM]}}
                {'hidden': {'team': ['round', 'time', 'racket', 'trunk', 'ketchup', 'millionaire', 'capital', 'circle', 'back'], 'innocent': ['eagle', 'olympus', 'lion'], 'opponent': ['worm', 'tooth', 'torch', 'board', 'spy', 'ground', 'card', 'watch'], 'assassin': ['suit', 'paper', 'post', 'missile', 'pipe']}, 'revealed': {'team': {'team': [], 'innocent': [], 'opponent': [], 'assassin': []}, 'opponent': {'team': [], 'innocent': [], 'opponent': [], 'assassin': []}}}
            }
        }
    }
    & & \\ \\

    \theutterance \stepcounter{utterance}  
    & & \multicolumn{4}{p{0.6\linewidth}}{
        \cellcolor[rgb]{0.9,0.9,0.9}{
            \makecell[{{p{\linewidth}}}]{
                \texttt{\tiny{[P1$\langle$GM]}}
                \texttt{You have attempted this same task before and did not succeed. These are the plans you wrote after those attempts.} \\
\texttt{Trial \#1: New plan: The previous attempt was aborted because the format included a newline between the CLUE and TARGETS lines, violating the "single{-}word clue" and "never include any other text" constraint which likely required a single{-}line output. Additionally, while "circle" is a valid clue for "round", I should ensure the TARGETS list is strictly comma{-}separated without extra spaces or line breaks. Next time, I will output the entire response on a single line in the exact format: "CLUE: <WORD> TARGETS: <WORD1>, <WORD2>, ...". I will also verify that the clue word is not one of the target words themselves (e.g., using "geometry" for "round" and "circle" instead of "circle" as the clue).} \\
\\ 
\\ 
\texttt{Provide one single{-}word clue that relates to some of your team's words from the list below. You can choose to how many and to which words your clue relates to.} \\
\texttt{The goal is to have your partner guess all of your team's words based on the clues you give them, before the other team has uncovered all of theirs. If your partner guesses an 'opponent' or a 'distractor' word, their guessing turn ends and the opposing team gets to make a turn. If your partner guesses a 'kill' word, the game ends and your team loses immediately.} \\
\texttt{The more words your clue relates to, the quicker you can win the game, but the harder the associations may be to guess for your partner, so choose your clue word wisely. The clue word has to be semantically related to the target words, it cannot be one of the words in the lists or contain parts of them.} \\
\\ 
\texttt{Always give your single{-}word clue and your comma{-}separated list of related target words in the following format and make your answers as short as possible, never include any other text than is required in this form:} \\
\\ 
\texttt{CLUE: <WORD>} \\
\texttt{TARGETS: <TARGETS>} \\
\\ 
\texttt{Your team words are: round, time, racket, trunk, ketchup, millionaire, capital, circle, back.} \\
\texttt{Your opponent's team words are: worm, tooth, torch, board, spy, ground, card, watch.} \\
\texttt{Distractor words are: eagle, olympus, lion.} \\
\texttt{The kill words are: suit, paper, post, missile, pipe.} \\
            }
        }
    }
    & \\ \\

    \theutterance \stepcounter{utterance}  
    & \multicolumn{4}{p{0.6\linewidth}}{
        \cellcolor[rgb]{0.9,0.9,0.9}{
            \makecell[{{p{\linewidth}}}]{
                \texttt{\tiny{[P1$\rangle$GM]}}
                \texttt{CLUE: geometry TARGETS: round, circle, back, capital} \\
            }
        }
    }
    & & \\ \\

    \theutterance \stepcounter{utterance}  
    & & & \multicolumn{2}{p{0.3\linewidth}}{
        \cellcolor[rgb]{0.9,0.9,0.9}{
            \makecell[{{p{\linewidth}}}]{
                \texttt{\tiny{[GM$|$GM]}}
                {'player': 'Cluegiver', 'type': 'prefix error', 'utterance': 'CLUE: geometry TARGETS: round, circle, back, capital', 'message': 'Your targets did not start with the correct prefix (TARGETS: ).', 'prefix': 'TARGETS: '}
            }
        }
    }
    & & \\ \\

    \theutterance \stepcounter{utterance}  
    & \multicolumn{4}{p{0.6\linewidth}}{
        \cellcolor[rgb]{0.95,0.95,0.95}{
            \makecell[{{p{\linewidth}}}]{
                \texttt{\tiny{[P2$\rightarrow$P1]}}
                \texttt{Post{-}episode Reflexion (not sent to GM):} \\
\texttt{New plan: The previous attempt was aborted because the output contained a newline character between the "CLUE" and "TARGETS" lines, violating the strict "single{-}line" format requirement. Additionally, the clue "geometry" is a valid semantic link for "round" and "circle", but I must ensure the entire response is one continuous string without any line breaks. Next time, I will construct the response as a single line: "CLUE: <WORD> TARGETS: <WORD1>, <WORD2>, ..." and verify that the clue word is not present in the target list or the forbidden word lists (opponent, distractor, kill) before sending. For example, using "shape" for "round" and "circle".} \\
            }
        }
    }
    & & \\ \\

\end{supertabular}
}

\end{document}
